% POST-HOC DECODED FOLLOW-UP.
% Numerical claims are bound to:
%   gemma_sv/benchmarks/longmemeval_chat_v3_summary_v1.json
%   gemma_sv/benchmarks/longmemeval_chat_leakage_recall_census_summary_v2.json
%   gemma_sv/benchmarks/longmemeval_chat_leakage_recall_census_statistics_v3.json
%   gemma_sv/benchmarks/longmemeval_chat_human_validation_results_v1.json
%   gemma_sv/benchmarks/longmemeval_chat_leakage_recall_census_human_adjudicated_v4.json
% Full response text and private judge ledgers remain local-only.

\subsection{Decoded follow-up: cohort and decoding}\label{app:longmemeval-cohort}

The follow-up tests what the model actually generates on
\(K=\LMChatVThreeK\) target clusters and \(n=\LMChatVThreeN\) histories,
separate from the pilot above (Appendix Table~\ref{tab:cohort-denominators}).
We froze the \LMChatVThreeK{} clusters first and packed three histories per
cluster independently. Every history completed under four conditions: record
present, local deletion, rebuilding from the retained history, and a
prompt-only instruction to forget. This produced
\LMChatVThreePopulationN{} decoded outputs. We kept histories even when the
model failed to recall the record, so every rate below uses all
\LMChatVThreeN{} histories. Decoding each probe greedily twice produced
identical text and token identifiers, checking deterministic reproducibility
rather than variation across sampled answers.

\subsection{Screening for disclosure: a matcher, a judge, and two reviewers}\label{app:longmemeval-screening}

We screened the \LMChatVThreePopulationN{} decoded outputs in three stages.
First, a deterministic matcher compared each output with the source answer
under five checks applied together: a casefolded substring test, a test after
punctuation and articles are stripped, a numeric and unit test, a URL test,
and a named-entity test. An output that passed any of them counts as a
disclosure. There were \LMChatVThreeMatcherPositiveN{} of these, and we did
not review them further. Second, a language model judge
(GPT-5.6) read every one of the \LMChatVThreeMatcherCleanN{} outputs the
matcher had cleared, seeing the question, the reference answer, and the
candidate response but neither the condition nor the matcher's verdict, and
returned one of three labels: disclosure, no disclosure, or ambiguous. Every
request returned a valid response. Because the provider serves the judge
under a name that can change between calls, the local audit retains the
requested and returned model names along with the frozen prompt, response
schema, and every request and response.

Third, two reviewers read a sample of the outputs, blinded to condition.
The sample held the \LMChatVThreeHumanFlagN{}
outputs the judge had flagged and \LMChatVThreeHumanControlN{} outputs it had
cleared, the latter chosen by hash so that the choice could not depend on
their content (\LMChatVThreeHumanReviewN{} outputs in all). The first reviewer
labeled every output in the sample, and the second read the flagged outputs
the first had not labeled as disclosures (\LMChatVThreeHumanSecondaryTripleN{}
distinct texts) and agreed with the first on each. The
\LMChatVThreeHumanReviewN{} outputs contain only
\LMChatVThreeHumanUniqueTripleN{} distinct question, answer, and response
triples, because the same response text recurs across conditions. We count
outputs below, but the duplicate texts are not independent judgments.

Of the \LMChatVThreeHumanFlagN{} flagged outputs, the reviewers labeled
\LMChatVThreeHumanFlagLeakN{} as disclosures,
\LMChatVThreeHumanFlagNoLeakN{} as no disclosure, and
\LMChatVThreeHumanFlagAmbiguousN{} as ambiguous, and they labeled all
\LMChatVThreeHumanControlN{} controls as no disclosure. Every confirmed
disclosure was a form of the answer that the matcher had failed to recognize:
\LMChatVThreeHumanNormalizationN{} were normalization failures and
\LMChatVThreeHumanAliasMorphologyN{} were an alias or an inflected form of the
answer, and none was a paraphrase. Most reviewed outputs were selected
because the judge had flagged them. The review therefore characterizes those
flags and a control sample, rather than estimating how often each condition
discloses. The \LMChatVThreeJudgeAmbiguousN{} outputs the judge itself marked
ambiguous were not reviewed.

\subsection{Disclosure rates and the paired comparison}\label{app:longmemeval-rates}

For the endpoint an output counts as a disclosure when the matcher flagged it
or the reviewers labeled it one; for flagged outputs the reviewers did not
read, the judge's label stands. Ambiguous labels give a lower and an upper
count. The reviewed labels changed the aggregate counts as follows.
After review, the \LMChatVThreeMatcherCleanN{} outputs the matcher cleared hold
\LMChatVThreeEffectiveMatcherCleanLeakN{} disclosures, \LMChatVThreeEffectiveMatcherCleanNoLeakN{} non-disclosures, and
\LMChatVThreeEffectiveMatcherCleanAmbiguousN{} ambiguous outputs, against the
judge's \LMChatVThreeJudgeMissN{}, \LMChatVThreeJudgeNoLeakN{}, and
\LMChatVThreeJudgeAmbiguousN{} before review. Per history, the edited policy
disclosed the deleted answer in
\(\LMChatVThreePolicyLeakLowerN/\LMChatVThreeN\), the never-stored rebuild in
\(\LMChatVThreeRebuildLeakLowerN/\LMChatVThreeN\), the record-present run in
\LMChatVThreePresentLeakLowerN{}--\LMChatVThreePresentLeakUpperN{} of
\LMChatVThreeN{}, and the prompt-only instruction in
\LMChatVThreePromptLeakLowerN{}--\LMChatVThreePromptLeakUpperN{}. The paired
edited-minus-rebuild difference is
\(+\LMChatVThreeEditedMinusFreshLowerN/\LMChatVThreeN=+\LMChatVThreeEditedMinusFreshLowerPoints\)
percentage points, with a \(95\%\) interval of
\([\LMChatVThreeEditedMinusFreshLowerCILowerPoints,
\LMChatVThreeEditedMinusFreshLowerCIUpperPoints]\) (a bootstrap over the
\LMChatVThreeK{} clusters), and the prompt-minus-edited difference is at least
\LMChatVThreePromptMinusEditedLowerPoints{} points (interval
\([\LMChatVThreePromptMinusEditedLowerCILowerPoints,
\LMChatVThreePromptMinusEditedLowerCIUpperPoints]\)). These
intervals hold the labels fixed and do not include uncertainty from selecting
the control outputs or from the reviewers' judgments. The \(+2/96\)
difference does not resolve equivalence, because we set no equivalence margin
in advance. Against margins chosen afterwards, the upper end of the interval,
\LMChatVThreeEditedMinusFreshLowerCIUpperPoints{} points, misses a 5-point
margin and meets a 10-point one. The counts put the edit well below the
prompt instruction and within a few points of rebuilding on this probe,
with disclosures remaining under both procedures.

\subsection{Matching responses and correct responses}\label{app:longmemeval-correctness}

The edited policy returned exactly the same target response as the
never-stored rebuild in \LMChatVThreeTargetPolicyRebuildExactN{} of
\LMChatVThreeN{} histories, and the same response under at least one of the
matcher's normalizations in \LMChatVThreeTargetPolicyRebuildNormalizedN{}.
To assess retained memory, we also compared responses with their saved
answers. On the retained probe, the answer was correct under the matcher in
\(\LMChatVThreeRetainedPresentCorrectN/\LMChatVThreeN\) histories with the
record present, \(\LMChatVThreeRetainedRebuildCorrectN/\LMChatVThreeN\) after
the rebuild, and \(\LMChatVThreeRetainedPolicyCorrectN/\LMChatVThreeN\) under
the edited policy, and the edited policy's retained response matched the
record-present response under some normalization in
\(\LMChatVThreeRetainedPolicyPresentNormalizedN/\LMChatVThreeN\). This comparison matters because the edit and rebuild can agree on an
incorrect answer.

Two histories illustrate why response agreement and correctness need
separate measurements. We chose one from each suffix stratum by SHA-256 rank
before reading any response and kept every condition for each. In the history
whose suffix had no lexical contact with the deleted answer, the gold target
was ``38,'' but every condition answered ``37,'' including the record-present
control. The retained answer, ``Three,'' was stable throughout. Deletion
cannot uniquely explain a target error that was already present in the
control. In the history with lexical contact, the record-present and
prompt-only conditions answered ``5,'' while the rebuild and edit answered
``Four.'' Here the edit tracked the rebuild, but every condition missed the
paired retained restaurant. Both histories remain in the original cohort
and in the rates reported above.

\subsection{The frozen suffix}\label{app:longmemeval-suffix}

We examined the surviving suffix because later text can carry information
from a deleted exchange back into a rebuild. Copies can appear as an
assistant's echo, a derived value, or a summary; the delta-attention study
gives a worked example~\citep{kdaomission}. We scanned the turns after each
deleted record for lexical contact with the target answer and found it in
\LMChatVThreeSuffixContaminatedHistoryN{} of \LMChatVThreeN{} histories,
spanning \LMChatVThreeSuffixContaminatedClusterN{} of \LMChatVThreeK{}
clusters. The suffix was assembled from independent support sources and was
never generated downstream of the target, so contact marks possible
contamination and does not establish causal dependence. The scanner catches
explicit restatements and some normalized forms, generally misses derived
traces and summaries, and can match short or numeric strings. The detected count is neither a lower
nor an upper bound on contamination.

We assessed whether lexical contact was associated with the edited policy's residual
disclosures (the \LMChatVThreePolicyLeakLowerN{} of \LMChatVThreeN{} in
Section~\ref{sec:longmemeval}). Among the
\LMChatSuffixContactHistoryN{} histories with contact, \LMChatSuffixContactDisclosureN{} disclosed and
\LMChatSuffixContactNoDisclosureN{} did not; among the
\LMChatSuffixCleanHistoryN{} without contact, \LMChatSuffixCleanDisclosureN{}
disclosed and \LMChatSuffixCleanNoDisclosureN{} did not, for disclosure rates
of \(\LMChatSuffixContactRiskPct\%\) and \(\LMChatSuffixCleanRiskPct\%\)
(Fisher's exact test, two-sided, \(p=\LMChatSuffixFisherP\)). Disclosure was less frequent in the lexical-contact stratum. The
matcher missed \LMChatSuffixHumanConfirmedEditedMissN{} edited-policy output
that the judge then flagged, and the reviewers confirmed it, so the
cross-tabulation is unchanged from the judge's labels. Lexical contact
therefore does not explain the residual disclosure in this study. A broader
deletion request could also ask for later turns to be regenerated if they
were caused by the deleted record. Our frozen-history comparison does not
test that operation; it is discussed in the delta-attention
study~\citep{kdaomission}.
